%==============================================================================
% jarvis-preamble.tex
% Shared packages, metadata macros, and semantic text commands for the
% Jarvis-HEP manual.  Loaded by main.tex after \documentclass{tufte-book}.
%
% The tufte-book class already provides: xcolor, hyperref (unless nohyper),
% geometry, graphicx, \allcaps, \smallcaps, \newthought, \sidenote,
% \marginnote, \footnote->sidenote, fullwidth env, marginfigure, etc.
% We only add what a code/reference manual additionally needs.
%==============================================================================

% --- Math ---------------------------------------------------------------------
\usepackage{amsmath}
\usepackage{amssymb}

% --- Tables -------------------------------------------------------------------
\usepackage{booktabs}
\usepackage{longtable}
\usepackage{array}
\usepackage{ragged2e}      % \RaggedRight inside narrow table cells

% --- Misc ---------------------------------------------------------------------
\usepackage{pmboxdraw}     % render block/box-drawing glyphs (the Jarvis -v banner) in pdfLaTeX
\usepackage{newunicodechar} % map the literal dot glyph U+2B24 (the banner icon) to a drawn circle
\usepackage{varwidth}      % shrink-to-content box so the banner panel hugs its text (no empty margin)
\usepackage{xspace}        % smart trailing space for text macros
\usepackage{enumitem}      % tighter description/itemize lists
\setlist{noitemsep, topsep=2pt}
% Bera Mono (the typewriter font tufte-common.def loads) carries the T1 `--' ->
% en-dash ligature, so every long CLI flag mentioned in prose via \cli{}/\code{}
% printed as `-strict' instead of `--strict'. microtype is loaded ONLY to switch
% that ligature off: protrusion and expansion are disabled so it changes nothing
% else about the book's typesetting (verified: page count identical, no new
% Overfull >15pt).
\usepackage[protrusion=false,expansion=false]{microtype}
\DisableLigatures[-]{family=tt*}

% --- Project version single-source --------------------------------------------
% This manual is a LIVING document: it is revised in step with each Jarvis-HEP
% release.  \JversionNum names the release this particular build documents --
% bump it (and rebuild) on every release; do not hard-code version numbers in
% the prose, always use \Jversion so there is a single source of truth.
\newcommand{\JversionNum}{2.0.0}
\newcommand{\Jversion}{v\JversionNum\xspace}

% --- Product / ecosystem names ------------------------------------------------
\newcommand{\Jarvis}{\textsc{Jarvis-HEP}\xspace}
\newcommand{\Jtwo}{\textsc{Jarvis-HEP}~2\xspace}
\newcommand{\Operas}{\textsc{Jarvis-Operas}\xspace}
\newcommand{\JPlot}{\textsc{JarvisPLOT}\xspace}
\newcommand{\Portal}{\textsc{Jarvis-HEP-Portal}\xspace}

% --- Inline semantic markup ---------------------------------------------------
% Monospace helper that lets long tokens break in the narrow text column: an
% underscore (\_) and a dot offer a zero-penalty breakpoint, so names like
% interpolations\_1D or Runtime.batch\_size can wrap instead of overflowing.
\newcommand{\jbttbreak}{\penalty0\hskip0pt\relax}
\newcommand{\jtt}[1]{\begingroup
  \def\_{\textunderscore\jbttbreak}%
  \texttt{#1}\endgroup}
% Code-ish things rendered in monospace with consistent meaning across the book.
\newcommand{\code}[1]{\jtt{#1}}                          % generic inline code
% A YAML key / block name -- coloured jYKey (style/jarvis-listings.tex) so it
% matches how the same key looks inside a yamlcode listing; \ykey{} is a
% shorter alias defined alongside that colour.
\newcommand{\yamlkey}[1]{\textcolor{jYKey}{\jtt{#1}}}
\newcommand{\cli}[1]{\jtt{#1}}   % a shell / CLI fragment -- real look defined in
                                  % jarvis-boxes.tex (\hl{} needs jShell + soul)
% \clih{}: plain fallback for \section{}/\chapter{}/\caption{} (moving arguments)
% -- soul's \hl{} cannot appear there (crashes: "Argument of \let has an extra }"),
% so any \cli{} inside a heading or caption must use \clih{} instead.
\newcommand{\clih}[1]{\jtt{#1}}
% \token{}/\marker{}: runtime tokens (@SampleID) and path markers (&J/) -- coloured
% jYValue (style/jarvis-listings.tex) since these are values substituted at run time,
% same colour family as a YAML string value in a listing.
\newcommand{\token}[1]{\textcolor{jYValue}{\jtt{#1}}}
\newcommand{\marker}[1]{\textcolor{jYValue}{\jtt{#1}}}
\newcommand{\file}[1]{\jtt{#1}}                          % a file or path (overridden in dirtree)
% \sampler{}: a sampler method name -- bold + jSampler purple, sans-serif kept.
\newcommand{\sampler}[1]{\textbf{\textcolor{jSampler}{\textsf{#1}}}}
\newcommand{\optitem}[1]{\textit{#1}}                    % an "(optional)" tag etc.

% Reference to a sampler JSON schema file (authoritative key source).
\newcommand{\schemaref}[1]{\texttt{card/schema/#1}}

% A "registered alias" inline note, e.g. \alias{ToyMCMC}{MCMC}
\newcommand{\alias}[2]{\texttt{#1}~(alias of~\texttt{#2})}

% --- Required / Optional inline tags for key tables ---------------------------
% Rendered as small inline badges (tcolorbox's \tcbox, "on line" keeps them in the
% running text/table cell without breaking it) -- colours jWarn (red) / jOpt (gold),
% both in style/jarvis-boxes.tex, which loads after this file but only needs to
% exist by the time \req/\opt are actually used in a chapter, not at definition time.
% The trailing \allowbreak{} gives TeX a legal line-break point right after the
% badge -- without it, a badge glued directly to following punctuation with no
% source space (e.g. "\req.") is one unbreakable unit with the next word and can
% force an overfull line; plain text never had this problem since a period alone
% is breakable. (A zero-width \hskip was tried first and did not fix it --
% \allowbreak's bare penalty is what actually works; keep the empty braces, a
% bare \allowbreak before a letter gets swallowed into an undefined control
% sequence.)
\newcommand{\req}{\tcbox[on line, colback=jWarn!10, colframe=jWarn, coltext=jWarn,
  boxrule=0.6pt, arc=1.5pt, left=3pt, right=3pt, top=0pt, bottom=0pt,
  boxsep=1pt]{\sffamily\bfseries\footnotesize REQ}\allowbreak{}}
\newcommand{\opt}{\tcbox[on line, colback=jOpt!18, colframe=jOpt, coltext=jOpt,
  boxrule=0.6pt, arc=1.5pt, left=3pt, right=3pt, top=0pt, bottom=0pt,
  boxsep=1pt]{\sffamily\bfseries\footnotesize OPT}\allowbreak{}}

% --- Tufte: \paragraph is our de-facto third heading level --------------------
% tufte-book forbids \subsubsection, so \paragraph (provided by the class) is the
% deepest heading we use.  We deliberately do NOT redefine it, to avoid clashing
% with the class's titlesec-based sectioning.

% --- A compact key/value description list for YAML key references -------------
\newenvironment{keylist}
  {\begin{description}[style=nextline, leftmargin=1.2em, font=\ttfamily]}
  {\end{description}}

% --- Coloured cross-reference & citation links --------------------------------
% tufte loads hyperref with colorlinks=false (links are black, box-less), so
% references look like plain text.  Turn on coloured links so references to
% figures, tables, equations, and sections -- and citations -- stand out.
% (hyperref is already loaded by the tufte-book class at this point.)
\definecolor{jlink}{HTML}{1F6FB2}  % cross-refs: figures, tables, equations, sections
\definecolor{jcite}{HTML}{2E7D32}  % citations / bibliography
\definecolor{jurl}{HTML}{1A5FB4}   % URLs
\hypersetup{
  colorlinks = true,
  linkcolor  = jlink,
  citecolor  = jcite,
  urlcolor   = jurl,
  linktoc    = page,   % keep ToC titles black; colour only the page-number links
}
